% GENERATED FILE. Edit canonical lessons, then run npm run generate:books.
% canonical-source-hash: fnv1a64:fd61b000316a040e
% canonical-lessons: KA-C235-bilu, KA-C235-hindirugu, KA-C235-muttu, KA-C235-ettu, KA-C235-ese

\chapter{To fall, To return, To touch, To lift, To throw}
\label{ch:ka-to-fall-to-return-to-touch-to-lift-to-throw}
\hlchaptermodality{235}

\begin{chapteropening}
\textbf{By the end of this chapter:} \emph{I can say to fall, to return, to touch, to lift and to throw.}

Complete the last lesson of chapter 235: I can say to fall, to return, to touch, to lift and to throw.
\end{chapteropening}

\section[\texorpdfstring{\kn{ಬೀಳು} (bīḷu)}{ಬೀಳು (bīḷu)}]{\kn{ಬೀಳು} (bīḷu) — to fall}

\label{lesson:KA-C235-bilu}



\begin{quote}
Before the new one: say the Kannada for a work shift, then the Kannada for an advance.
\end{quote}

\subsection*{You'll want to know: \kn{ಬೀಳು}}
\textbf{\kn{ಬೀಳು}} — \emph{bīḷu} — \textquotedblleft{}to fall\textquotedblright{}.

\begin{grammarlens}[title={What you've built}]
One more word for this chapter.
\end{grammarlens}

\subsection*{Your turn}
\begin{itemize}
\raggedright
  \item \emph{Say it:} \emph{bīḷu}
  \item \emph{Say it:} \emph{bīḷu}, once more
  \item \emph{Say it:} say \emph{bīḷu}
  \item \emph{Recall:} say the Kannada for a work shift, then the Kannada for an advance, then say \emph{bīḷu} again
\end{itemize}

\begin{tcolorbox}[breakable,colback=teal!4,colframe=teal!35!black,title={Before you move on}]
What does \kn{ಬೀಳು} mean? (\textquotedblleft{}To fall\textquotedblright{}.) Say it once more.
\end{tcolorbox}
\section[\texorpdfstring{\kn{ಹಿಂದಿರುಗು} (hindirugu)}{ಹಿಂದಿರುಗು (hindirugu)}]{\kn{ಹಿಂದಿರುಗು} (hindirugu) — to return, to come back}

\label{lesson:KA-C235-hindirugu}



\begin{quote}
Before the new one: say the Kannada for an advance, then the Kannada for to fall.
\end{quote}

\subsection*{You'll want to know: \kn{ಹಿಂದಿರುಗು}}
\textbf{\kn{ಹಿಂದಿರುಗು}} — \emph{hindirugu} — \textquotedblleft{}to return, to come back\textquotedblright{}.

\begin{grammarlens}[title={What you've built}]
One more word for this chapter.
\end{grammarlens}

\subsection*{Your turn}
\begin{itemize}
\raggedright
  \item \emph{Say it:} \emph{hindirugu}
  \item \emph{Say it:} \emph{hindirugu}, once more
  \item \emph{Say it:} say \emph{hindirugu}
  \item \emph{Recall:} say the Kannada for an advance, then the Kannada for to fall, then say \emph{hindirugu} again
\end{itemize}

\begin{tcolorbox}[breakable,colback=teal!4,colframe=teal!35!black,title={Before you move on}]
What does \kn{ಹಿಂದಿರುಗು} mean? (\textquotedblleft{}To return, to come back\textquotedblright{}.) Say it once more.
\end{tcolorbox}
\section[\texorpdfstring{\kn{ಮುಟ್ಟು} (muṭṭu)}{ಮುಟ್ಟು (muṭṭu)}]{\kn{ಮುಟ್ಟು} (muṭṭu) — to touch}

\label{lesson:KA-C235-muttu}



\begin{quote}
Before the new one: say the Kannada for to fall, then the Kannada for to return.
\end{quote}

\subsection*{You'll want to know: \kn{ಮುಟ್ಟು}}
\textbf{\kn{ಮುಟ್ಟು}} — \emph{muṭṭu} — \textquotedblleft{}to touch\textquotedblright{}.

\begin{grammarlens}[title={What you've built}]
One more word for this chapter.
\end{grammarlens}

\subsection*{Your turn}
\begin{itemize}
\raggedright
  \item \emph{Say it:} \emph{muṭṭu}
  \item \emph{Say it:} \emph{muṭṭu}, once more
  \item \emph{Say it:} say \emph{muṭṭu}
  \item \emph{Recall:} say the Kannada for to fall, then the Kannada for to return, then say \emph{muṭṭu} again
\end{itemize}

\begin{tcolorbox}[breakable,colback=teal!4,colframe=teal!35!black,title={Before you move on}]
What does \kn{ಮುಟ್ಟು} mean? (\textquotedblleft{}To touch\textquotedblright{}.) Say it once more.
\end{tcolorbox}
\section[\texorpdfstring{\kn{ಎತ್ತು} (ettu)}{ಎತ್ತು (ettu)}]{\kn{ಎತ್ತು} (ettu) — to lift}

\label{lesson:KA-C235-ettu}



\begin{quote}
Before the new one: say the Kannada for to return, then the Kannada for to touch.
\end{quote}

\subsection*{You'll want to know: \kn{ಎತ್ತು}}
\textbf{\kn{ಎತ್ತು}} — \emph{ettu} — \textquotedblleft{}to lift\textquotedblright{}.

\begin{grammarlens}[title={What you've built}]
One more word for this chapter.
\end{grammarlens}

\subsection*{Your turn}
\begin{itemize}
\raggedright
  \item \emph{Say it:} \emph{ettu}
  \item \emph{Say it:} \emph{ettu}, once more
  \item \emph{Say it:} say \emph{ettu}
  \item \emph{Recall:} say the Kannada for to return, then the Kannada for to touch, then say \emph{ettu} again
\end{itemize}

\begin{tcolorbox}[breakable,colback=teal!4,colframe=teal!35!black,title={Before you move on}]
What does \kn{ಎತ್ತು} mean? (\textquotedblleft{}To lift\textquotedblright{}.) Say it once more.
\end{tcolorbox}
\section[\texorpdfstring{\kn{ಎಸೆ} (ese)}{ಎಸೆ (ese)}]{\kn{ಎಸೆ} (ese) — to throw}

\label{lesson:KA-C235-ese}



\begin{quote}
Before the new one: say the Kannada for to touch, then the Kannada for to lift.
\end{quote}

\subsection*{You'll want to know: \kn{ಎಸೆ}}
\textbf{\kn{ಎಸೆ}} — \emph{ese} — \textquotedblleft{}to throw\textquotedblright{}.

\begin{grammarlens}[title={What you've built}]
One more word for this chapter.
\end{grammarlens}

\subsection*{Your turn}
\begin{itemize}
\raggedright
  \item \emph{Say it:} \emph{ese}
  \item \emph{Say it:} \emph{ese}, once more
  \item \emph{Say it:} say \emph{ese}
  \item \emph{Recall:} say the Kannada for to touch, then the Kannada for to lift, then say \emph{ese} again
\end{itemize}

\begin{tcolorbox}[breakable,colback=teal!4,colframe=teal!35!black,title={Before you move on}]
What does \kn{ಎಸೆ} mean? (\textquotedblleft{}To throw\textquotedblright{}.) Say it once more.
\end{tcolorbox}
